\chapter{PINN Benchmark Protocol}
\label{app:pinn_pricing_benchmark_setup}

This appendix records the evaluation protocol used for the PINN results in Chapter~\ref{ch:performance_benchmarks}. The methodological definitions of the baseline PINN, the Sobolev extension, and the analytical control variate are given in Chapter~\ref{ch:pinn_greeks}; here we only fix the conventions needed to reproduce the reported benchmarks.

\section{Evaluated Runs}

The final PINN comparison uses three stored runs:
\begin{table}[H]
    \ttabbox[\FBwidth]
    {\caption{Final PINN artefacts used in the benchmark}\label{tab:pinn_final_artifacts}}
    {\begin{tabular}{|P{3.0cm}|P{4.2cm}|P{4.2cm}|}
        \hline
        \textbf{Variant} & \textbf{Directory} & \textbf{Run name} \\
        \hline
        PINN baseline & \texttt{outputs/pinn} & \texttt{PINN\_mix\_scaled}\newline\texttt{fixed\_theta} \\
        \hline
        PINN Sobolev & \texttt{outputs/protos}\newline\texttt{sobolev\_mix}\newline\texttt{fixed\_theta} & \texttt{PINN\_mix\_scaled}\newline\texttt{fixed\_theta}\newline\texttt{20260420\_225517} \\
        \hline
        PINN Sobolev+ACV & \texttt{outputs/pinn} & \texttt{acv\_hard\_patch}\newline\texttt{sobolev\_control}\newline\texttt{variate\_best\_gate}\newline\texttt{tau\_floor\_5e4} \\
        \hline
        \multicolumn{3}{p{11.4cm}}{Source: Author's elaboration from stored experiment artefacts}
    \end{tabular}}
\end{table}
These artefacts are evaluated without retraining. The baseline and Sobolev PINNs are compared on the common full-grid protocol, while the Sobolev+ACV variant is used for the localized hard-region diagnostic defined in Section~\ref{sec:pinn_acv_correction}.

\section{Input And Reference Conventions}

All PINN evaluations use the 8-dimensional input convention
\[
\left[\tau, m, v, \rho, \kappa, \gamma, \bar v, r\right],
\]
where \(v\) is the current variance state. When the benchmark is evaluated at the calibrated market configuration, this coordinate is fixed to the calibrated value \(v_0^*\), while \((\rho^*,\kappa^*,\gamma^*,\bar v^*)\) define the Heston operator used by the trained model. Each input feature \(x_j\) is standardized with the affine map
\[
\tilde{x}_j=\frac{x_j-\mu_j^{\mathrm{train}}}{\sigma_j^{\mathrm{train}}},
\]
using the training statistics stored with each run. These statistics are not refitted on the benchmark grids.

The numerical reference is the Heston--COS engine introduced in Section~\ref{sec:fourier_cos_method}, under the European put convention with normalized strike \(K=1\). The COS configuration uses \(N=1500\), \(L=50\), and fallback \(L=80\) in the same stability cases described in Appendix~\ref{app:cos_error}. Reference Greeks are computed from the Heston characteristic-function formulas of Section~\ref{sec:heston_cf_greeks}. All reported benchmark quantities are evaluated in double precision on CPU.

\section{Benchmark Grids And Metrics}

The full-grid PINN benchmark contains \(25\,921\) valid points over
\[
m\in[0.8,1.2],\qquad \tau\in[0.05,2.0].
\]
This grid is used for the baseline and Sobolev pricing comparison and for the full-grid Greek ablation in Section~\ref{sec:pinn_greeks_performance}. Metrics follow the definitions in Section~\ref{sec:benchmark_error_metrics}; the appendix does not repeat the formulas, since the same MSE, RMSE, and MAPE conventions are used throughout Chapter~\ref{ch:performance_benchmarks}.

The Sobolev+ACV diagnostic is evaluated on the local short-maturity at-the-money region
\[
m\in[0.9704,1.0305],\qquad \tau\in[0.005,0.05].
\]
This separate grid is used only to test the hard-region correction. It is therefore not treated as a replacement for the full-grid Sobolev benchmark, but as a focused diagnostic for the remaining kink-driven Greek errors.
